\documentclass[11pt]{article}
\usepackage[a4paper,margin=26mm]{geometry}
\usepackage{amsmath,amssymb,amsthm}
\usepackage[hidelinks]{hyperref}
\newtheorem{theorem}{Theorem}
\title{A Prime Cyclic Cayley Graph:\\A Disproof of Conjecture 00000003964}
\author{GitHub: gaochengzhecpu}
\date{}
\begin{document}
\maketitle
\begin{abstract}
The source quantifies over every finite simple group, without excluding
cyclic groups of prime order. The Cayley graph of $\mathbb Z/1021\mathbb Z$
with steps $\pm1$ has diameter at least $510$, whereas the proposed upper
bound is less than $300$. The counterexample works with either natural
or binary logarithms. Its simplicity, generation, actual graph-distance
lower bound and logarithmic comparisons are verified in Lean.
\end{abstract}

We use the standard definition of a simple group: a nontrivial group
with no proper nontrivial normal subgroup. In particular, a cyclic group
of prime order is simple. Both languages of the source say ``every finite
simple group''; neither imposes a nonabelian hypothesis. For an additive
group generated by $S$, we use the undirected Cayley graph with edges
between distinct $u,v$ when $u-v\in S$ or $v-u\in S$. Thus its graph
distance is the word metric allowing generators and their inverses.

\begin{theorem}
Let $G=\mathbb Z/1021\mathbb Z$ and $S=\{1\}$. Then $G$ is a finite
simple group, $S$ generates $G$, and its Cayley graph $\Gamma$ satisfies
\[
 \begin{aligned}
 \operatorname{diam}\Gamma&\ge510>3(\ln |G|)^2,\\
 \operatorname{diam}\Gamma&>3(\log_2 |G|)^2.
 \end{aligned}
\]
\end{theorem}
\begin{proof}
The primes below $32$ are
\[
 2,3,5,7,11,13,17,19,23,29,31.
\]
None divides $1021$. Since $\sqrt{1021}<32$, the number $1021$ is prime.
Hence $G$ has prime order and is simple. Every residue is a nonnegative integer multiple
of $1$, so $S$ generates it. Its undirected Cayley graph is the cycle
on residues $0,1,\ldots,1020$.

For the representative $a\in\{0,\ldots,1020\}$ of a vertex, set
\[
 f(a)=\min\{a,1021-a\}.
\]
Along each edge the potential changes by at most one. Indeed consecutive
representatives change each of the two arguments of the minimum by at
most one; the wraparound edge joins $1020$ and $0$, whose potentials are
$1$ and $0$. For a walk of length $\ell$ from $u$ to $v$, induction gives
$f(v)\le f(u)+\ell$. Since $f(0)=0$ and $f(510)=510$, every walk from
$0$ to $510$ has length at least $510$. The graph is connected, so its
actual shortest-path distance, and hence its diameter, is at least $510$.

Finally, $1021<1024=2^{10}$ and $0<\ln2<1$ give
\[
 0<\ln1021<10\ln2<10,\qquad 0<\log_2 1021<10.
\]
Both proposed upper bounds are therefore strictly less than $300<510$.
\end{proof}

\section*{Formal verification and scope}
The Lean group is the actual type \texttt{ZMod 1021}. A kernel-checked
primality proof supplies the library's \texttt{IsSimpleAddGroup} instance.
The theorem \texttt{generator\_generates} establishes that the actual
additive-subgroup closure of $\{1\}$ is the whole group.

The graph is Mathlib's \texttt{cycleGraph 1021}, definitionally equal to
\texttt{circulantGraph (\{1\} : Set G)}. The library graph symmetrizes the
generating set. Connectivity is checked. The potential inequality is
proved for every actual edge, extended to every actual walk by induction,
and applied to a shortest walk supplied by the connected-graph API.
Thus \texttt{distance\_lower} concerns Mathlib's actual graph distance,
not a separately assigned numerical value.

The definition \texttt{diameter} is the maximum of these actual distances
over all vertex pairs. Lean proves that each distance is at most this
maximum, then obtains \texttt{diameter\_lower}. The theorem
\texttt{proposed\_bound\_lt\_diameter} checks the strict comparison with
the natural-logarithm expression. The additional theorem
\texttt{binary\_bound\_lt\_diameter} checks the quotient
$\ln|G|/\ln2$ as well. Finally, \texttt{conjecture\_false} negates the
universal diameter upper bound for finite simple groups, written in
additive notation, with a generating-set hypothesis and connectivity.

This refutes the source as stated. It does not address a different
conjecture restricted to nonabelian finite simple groups. No assertion
about the source's additional $\operatorname{PSL}(2,p)$ extremal-family
clause is needed, because its universal upper bound already fails.

\section*{Source and reproduction}
The exact bilingual source is included byte-for-byte as \texttt{SOURCE.md};
its repository path is \texttt{conjectures/00000003964.md}. The project
pins Lean 4.19.0 and Mathlib commit
\begin{center}\small\texttt{c44e0c8ee63ca166450922a373c7409c5d26b00b}.\end{center}
From \texttt{lean/}, run \texttt{lake build}, then
\begin{center}\small\texttt{lake env lean -DwarningAsError=true Main.lean}.\end{center}
The included public-Git manifest locks all dependency revisions.
The \texttt{verification/} directory records actual fresh-build logs,
theorem-axiom reports, artifact hashes, review and final PDF inspection.
No supplementary numerical computation is required.
\end{document}
